\section*{Overview}

Rollback DSU is the union-find structure I want when connectivity changes over time but the solution can be processed
offline. The ordinary DSU is great at merging components, but it cannot undo path compression. Rollback DSU fixes that
by giving up path compression and recording enough information to reverse each merge.

\section*{When to Use It}

Use rollback DSU when:

\begin{itemize}
  \item unions need to be undone later,
  \item the problem is offline, often over a segment tree on time,
  \item connectivity or component data changes across intervals of activity.
\end{itemize}

The standard use case is offline dynamic connectivity: edges appear and disappear over time, and queries ask whether
vertices are connected at specific moments.

\section*{Core Idea}

Keep union by size or rank, but avoid path compression. Every successful merge pushes enough information onto a stack to
restore the previous parents, sizes, and component count.

Then the structure supports:
\begin{itemize}
  \item \texttt{snapshot()}: remember the current stack size,
  \item \texttt{rollback(s)}: undo everything after that snapshot.
\end{itemize}

\section*{Key Insight}

Rollback only works cleanly because the parent pointers change in a controlled way. Path compression changes many nodes
at once and makes undo bookkeeping too expensive, so rollback DSU intentionally uses a slightly slower but reversible
\texttt{find}.

\section*{Operations / Main Technique}

\begin{itemize}
  \item \texttt{find} without compression,
  \item \texttt{unite} with union by size,
  \item store each successful merge on a history stack,
  \item restore by popping the stack.
\end{itemize}

\paragraph{Worked example.}
If I merge components of sizes \(3\) and \(5\), I only need to remember which root became the child and what the old
size of the larger root was. Undoing the merge restores those two values and increments the component count.

\section*{Correctness Intuition}

Each successful union changes only a constant amount of state: one parent pointer, one size entry, and the number of
components. Because those exact old values are stored, rollback restores the DSU to a previous valid state exactly.

\section*{Complexity Analysis}

\begin{itemize}
  \item \texttt{find}: \(O(\log n)\) amortized with union by size,
  \item \texttt{unite}: \(O(\log n)\),
  \item \texttt{rollback}: \(O(1)\) per undone merge,
  \item memory: \(O(n + \text{number of successful merges})\).
\end{itemize}

\section*{Implementation}

The code exposes:

\begin{itemize}
  \item \texttt{find},
  \item \texttt{unite},
  \item \texttt{same},
  \item \texttt{snapshot},
  \item \texttt{rollback}.
\end{itemize}

That is the core I would plug into a segment-tree-over-time solution.

\section*{Common Pitfalls}

\begin{itemize}
  \item Accidentally adding path compression and breaking rollback.
  \item Forgetting to record failed merges distinctly from successful ones.
  \item Rolling back to the wrong snapshot boundary inside recursion.
  \item Assuming rollback DSU is a drop-in replacement for online dynamic connectivity.
\end{itemize}

\section*{Variants / Extensions}

\begin{itemize}
  \item Component metadata such as sum, parity, or bipartite flags with rollback support.
  \item Segment tree over time for edge-activity intervals.
  \item Persistent DSU, which is a different design with stronger version access.
\end{itemize}

\section*{Practice Problems}

\begin{itemize}
  \item Offline dynamic connectivity with add, remove, and connectivity queries.
  \item Bipartite checking under interval-active edges.
  \item Time-segment recursion where each branch needs a reversible DSU state.
\end{itemize}

\section*{References}

\begin{itemize}
  \item \href{https://cp-algorithms.com/data_structures/deleting_in_log_n.html}{cp-algorithms: Deleting from a Data Structure in \(O(T(n)\log n)\)}.
  \item \href{https://codeforces.com/catalog?locale=en}{Codeforces Catalog}.
  \item \href{https://github.com/kth-competitive-programming/kactl}{KACTL}.
\end{itemize}
